%!TEX root = ../thesis.tex
%*******************************************************************************
%****************************** Chapter 2 **************************************
%*******************************************************************************

\chapter{PROJECT MANAGEMENT PLAN}
\label{chap:pm}

\ifpdf
    \graphicspath{{Chapter2/Figs/Raster/}{Chapter2/Figs/PDF/}{Chapter2/Figs/}}
\else
    \graphicspath{{Chapter2/Figs/Vector/}{Chapter2/Figs/}}
\fi

This chapter describes how the team organised the 15 weeks of the project, following the
structure of the project management plan in the AIP491 capstone guide.
Section~\ref{sec:intro-team} presents the team and how it communicates.
Section~\ref{sec:intro-pm} presents the project management approach: the phases, the work
breakdown and the schedule, and how risks, quality, the compute budget, changes and the
closure of the project were managed. Sections~\ref{sec:pm-training} to~\ref{sec:pm-config}
cover the training plan, the deliverables and configuration management. The project is
research, so the plan revolves around experiments: every conclusion in
Chapter~\ref{chap:results} must be traceable to a run with a log, a seed and a fixed
configuration.

%********************************** Section 2.1 ********************************
\section{Team work}
\label{sec:intro-team}

\subsection{Team structure and roles}

The team has four students, each in charge of one technical area as defined in the
project proposal.

Pham Le Huy Hoang (SE182670) is the team leader and is in charge of PEFT / Training.
Besides training the QLoRA models, running the LoRA configuration ablations and
building the shared adapter, Hoang coordinates the schedule, integrates the results
of the different areas and reviews cross-area results.

Ho Dac Toan (SE183249) is in charge of VLM / Baseline: studying the Qwen2.5-VL
architecture, reproducing the zero-shot baseline and comparing it with the published
numbers.

Tran Quoc Huan (SE193642) is in charge of Dataset / OCR: obtaining the data from the
official sources, fixing the splits, checking for leakage between training and test,
and building the OCR cache with PaddleOCR and VietOCR.

Van Quang Truong (SE182458) is in charge of Evaluation / Deployment: implementing the
metrics and statistical tests, measuring hallucination, and building the prototype and
the demo video.

Cross-area work, for example B3, which needs SV2's OCR cache and SV3's training pipeline, is
the responsibility of the person producing the main result and is reviewed by the team
leader. Each member's specific contributions by week are recorded in the Team Reports.
Table~\ref{tab:pm-roles} summarises the assignments, with each member's main source-code
directories.

\begin{table}[htbp]
\caption{Assignments}
\label{tab:pm-roles}
\centering
\small
\begin{tabular}{p{0.20\textwidth}p{0.18\textwidth}p{0.32\textwidth}p{0.20\textwidth}}
\toprule
\textbf{Member} & \textbf{Role} & \textbf{Main responsibilities} & \textbf{Main directories} \\
\midrule
Ho Dac Toan (SE183249) & SV1: VLM / Baseline & Understand the Qwen2.5-VL architecture, reproduce the zero-shot baseline, compare with published numbers & \path{src/models/}, \path{src/eval/run_baseline.py} \\
Tran Quoc Huan (SE193642) & SV2: Dataset / OCR & Official data, splits, leakage checks, OCR cache with PaddleOCR and VietOCR & \path{data/}, \path{src/data/} \\
Pham Le Huy Hoang (SE182670) & SV3: PEFT / Training; team leader & QLoRA training, LoRA configuration ablations, shared adapter; coordination, integration & \path{src/adaptation/}, \path{configs/}, \path{adapters/} \\
Van Quang Truong (SE182458) & SV4: Evaluation / Deployment & Metrics, statistical tests, hallucination, prototype and demo & \path{src/eval/}, \path{src/prototype/} \\
\bottomrule
\end{tabular}
\end{table}

\subsection{Communication plan}

The team works with the supervisors at the department's milestones and whenever a
major decision needs confirmation, for example fixing the Qwen2.5-VL-3B baseline on
7 October 2026. After each review meeting, the team fills in the feedback checklist,
emails it to the supervisors and the reviewers, and uploads it to the shared folder.

Within the team, each working week has a README and a Team Report in the
\path{docs/Wxx/} directory of the repository, recording what was done, the results,
the decisions and what was postponed. Source code, small logs and documents are
managed on GitHub; images, adapters, the OCR cache and the demo video are stored on
Google Drive. Decisions that affect the main results are recorded before the run and
confirmed by the supervisors when needed.

Table~\ref{tab:pm-comm} lists the communication channels and how often each is used.

\begin{table}[htbp]
\caption{Communication channels}
\label{tab:pm-comm}
\centering
\small
\begin{tabular}{p{0.26\textwidth}p{0.48\textwidth}p{0.16\textwidth}}
\toprule
\textbf{Channel} & \textbf{Purpose} & \textbf{Frequency} \\
\midrule
Meetings with the supervisors & Confirm major decisions, e.g.\ fixing the Qwen2.5-VL-3B baseline on 7 October 2026 & At milestones and when needed \\
Reviews with the reviewers & External feedback; after each review the team fills in the AI\_Capstone\_Review123 checklist, emails it to the supervisors and the reviewers, and uploads it to the shared folder & Twice \\
Weekly reports (README and Team Report in \path{docs/Wxx/}) & Record what was done, the results, the decisions and what was postponed & Every working week \\
GitHub & Source code, small logs, documents; every commit has a description & Continuous \\
Google Drive & Images, adapters, OCR cache, demo video, review recordings & Continuous \\
Department email & Review date announcements, sent only in review weeks & At milestones \\
\bottomrule
\end{tabular}
\end{table}

%********************************** Section 2.2 ********************************
\section{Project management approach}
\label{sec:intro-pm}

This section describes how the work was planned and controlled: the phases of the project,
the work breakdown and the schedule, and how risks, quality, the compute budget, changes and
the closure of the project were managed.

\subsection{Phases}

The project is divided into five phases following the 15-week frame of the proposal:

\paragraph{Phase 1: Orientation and planning (W01--W03).} The team reads papers in the
four areas, builds a literature matrix, analyses the two datasets, identifies the
research gaps, research questions and hypotheses, and presents them at Review~1.

\paragraph{Phase 2: Data and baseline (W04--W07).} The team obtains the data from the
official sources, fixes the Qwen2.5-VL-3B-Instruct baseline, measures zero-shot,
reproduces the published number on TextVQA and fixes the inference configuration
based on measurements (number of image tokens, data type).

\paragraph{Phase 3: Adaptation and improvement (W07--W10).} The team trains plain
QLoRA, analyses the errors, proposes OCR-aware QLoRA, verifies it on the
8,000-question subset, checks it across three seeds and with the $2\times2$ table, and
then runs it on the full data.

\paragraph{Phase 4: Evaluation and analysis (W10--W12).} The team builds the full
evaluation table, runs the LoRA configuration ablations, tries a shared adapter,
analyses the errors and measures hallucination with the Vietnamese POPE benchmark.

\paragraph{Phase 5: Prototype, packaging and thesis (W13--W15).} The team builds the web
prototype, verifies that it matches the evaluation pipeline, records the demo video,
and prepares the reproducibility package, the defence slides and the thesis.

\subsection{Work breakdown structure}
\label{sec:pm-wbs}

The work is divided into eight work packages (Table~\ref{tab:pm-wbs}). Time is counted in
working weeks W01--W15 within the 15-week frame of the proposal;
Section~\ref{sec:pm-process} shows how these weeks map onto the actual calendar.

\begin{table}[htbp]
\caption{Work breakdown}
\label{tab:pm-wbs}
\centering
\small
\begin{tabular}{p{0.22\textwidth}p{0.40\textwidth}p{0.13\textwidth}p{0.14\textwidth}}
\toprule
\textbf{Work package} & \textbf{Main content} & \textbf{Weeks} & \textbf{Main owner} \\
\midrule
1. Survey and orientation & Read papers by area, literature matrix, research gaps, research questions, hypotheses & W01--W03 & Whole team \\
2. Data & Obtain ViVQA and ViTextVQA from the official sources, fix splits, check train/test leakage, OCR cache & W02, W05, W08, W14 & SV2 \\
3. Baseline & Reproduce Qwen2.5-VL-3B-Instruct zero-shot (B1), compare with the published TextVQA number (R0), fix the inference configuration & W04--W07 & SV1 \\
4. Adaptation with QLoRA & QLoRA training pipeline; models B2, B3, B3-full, C3, shared adapter SH; LoRA configuration ablations & W06--W11 & SV3 \\
5. Evaluation & EM, VQA-Acc, ANLS, F1; paired McNemar tests; stability across seeds; full evaluation table & W05--W11 & SV4 \\
6. Analysis & Error analysis, question classification, POPE-vi hallucination, token budget, effect of normalisation and word segmentation on the metrics & W07, W11--W12, W15 & SV4, SV3 \\
7. Prototype & Web demo, yes/no threshold calibration, checking that the prototype matches the evaluated system, demo video & W13--W14 & SV4 \\
8. Reports and defence & Weekly reports, review slides, reproducibility package, thesis & Weekly, W15 & Whole team, coordinated by the leader \\
\bottomrule
\end{tabular}
\end{table}

\paragraph{Management objectives.} The plan is managed against four objectives:

\begin{itemize}
  \item \textbf{Reproducibility:} every run records split, seed, hardware, hyperparameters,
  training and inference time.
  \item \textbf{No choosing results after seeing the numbers:} the decision rule of every main
  test is committed to Git before running it.
  \item \textbf{Within the resource constraints:} QLoRA only, every experiment runs on one
  Colab GPU (T4 or L4).
  \item \textbf{On schedule:} the main research is finished before Review~2 and the
  department committee.
\end{itemize}

By the end of W14, the \path{experiments/} directory holds 141 runs, each with a log. The six
official models are listed in \path{configs/official_models.json}, a file generated
automatically from the logs; the four adapters stored on Google Drive all match the recorded
SHA-256 hashes.

\subsection{Schedule and progress tracking}
\label{sec:pm-process}

The team works in weekly cycles. Each working week goes through five steps:

\begin{enumerate}
  \item \textbf{Plan:} choose the week's questions from the overall plan and the previous
  week's results, and record them in the week's README.
  \item \textbf{Record rules in advance:} for each main test, write the decision rule
  (threshold, test, dataset) and commit it to Git before running.
  \item \textbf{Run:} code is written and self-tested on a personal machine, then pushed to
  GitHub; a Colab notebook clones the repository and runs on a GPU; logs and results are
  stored on Google Drive.
  \item \textbf{Check:} recompute the numbers from the original prediction files
  (\path{predictions.json}) instead of copying them by hand, and re-clone the repository to
  verify each commit.
  \item \textbf{Close the week:} update the week's README, the Team Report and the draft of
  the relevant chapters.
\end{enumerate}

The working weeks W01--W15 are numbered by content within the proposal's frame, not by the
calendar. The team ran ahead of schedule: on 2 November 2026 (calendar week 9), the team had
already reached W15. Calendar weeks are counted from Monday 7 September 2026 as week 1.
Table~\ref{tab:pm-weeks} maps the two counts.

\begin{table}[htbp]
\caption{Working weeks and the actual calendar (2026)}
\label{tab:pm-weeks}
\centering
\small
\begin{tabular}{p{0.12\textwidth}p{0.55\textwidth}p{0.22\textwidth}}
\toprule
\textbf{Week} & \textbf{Main content} & \textbf{Actual dates} \\
\midrule
W01--W02 & Orientation, 15-week plan, literature matrix, data analysis & 7--16 Sep \\
W03 & Review~1 package & 17 Sep (Review~1: 24 Sep) \\
W04 & First zero-shot baseline; baseline changed to Qwen2.5-VL-3B & 18--21 Sep \\
W05 & Official data; zero-shot baseline on both datasets & 22--27 Sep \\
W06 & fp16 bug diagnosis; QLoRA pipeline & 28--29 Sep \\
W07 & Configuration fixed on the L4; official B1; B2; R0 reproduction & 30 Sep--12 Oct \\
W08 & OCR-enhanced prompting; B3 on the subset & 12--14 Oct \\
W09 & Stability across seeds; $2\times2$ table; 2-epoch ViVQA & 15--17 Oct \\
W10 & LoRA configuration ablations; B3-full & 19--21 Oct \\
W11 & Full evaluation table; shared adapter & 21--22 Oct \\
W12 & POPE-vi hallucination; error analysis & 23--24 Oct \\
W13 & Web prototype & 26--31 Oct \\
W14 & Code finalised, demo video, reproducibility package & 1--2 Nov \\
W15 & Review~2 feedback, defence slides, thesis & from 2 Nov \\
\bottomrule
\end{tabular}
\end{table}

The department's milestones follow the actual calendar (Table~\ref{tab:pm-milestones}). The
thesis grade is 50\% process, for which the team submits 7 reports to the supervisors, and 50\% defence. The seven reports correspond to the chapters of this thesis as shown in Table~\ref{tab:intro-aip491}.

\begin{table}[htbp]
\caption{Department milestones}
\label{tab:pm-milestones}
\centering
\small
\begin{tabular}{p{0.18\textwidth}p{0.22\textwidth}p{0.34\textwidth}p{0.15\textwidth}}
\toprule
\textbf{Milestone} & \textbf{Date} & \textbf{Content} & \textbf{Status} \\
\midrule
Review~1 & 24 Sep 2026 (calendar week 3) & Problem, gaps, research questions, plan; feedback from two lecturers other than the supervisors & Done, 10 comments \\
Review~2 & 29 Oct 2026 (calendar week 8) & Reproduced baseline, method, preliminary results & Done \\
Department committee & 30 Nov--6 Dec 2026 (calendar week 13) & Three members, 60 minutes per team & Upcoming \\
University committee & 14--20 Dec 2026 (calendar week 15) & Five members & Upcoming \\
\bottomrule
\end{tabular}
\end{table}

\subsection{Risk management}
\label{sec:pm-risk}

The W01 plan listed eight risks (model too large, lack of GPUs, faulty data, baseline hard to
reproduce, poor OCR, training not converging, demo hard to deploy, slow training).
Table~\ref{tab:pm-risks} records the risks that actually occurred and how the team handled
them.

\begin{table}[htbp]
\caption{Risks that occurred and how they were handled}
\label{tab:pm-risks}
\centering
\small
\begin{tabular}{p{0.24\textwidth}p{0.08\textwidth}p{0.10\textwidth}p{0.46\textwidth}}
\toprule
\textbf{Risk} & \textbf{Week} & \textbf{Level} & \textbf{Handling} \\
\midrule
No university GPU server & W06 & High & Moved all training to Colab; upgraded to Colab Pro on 7 October 2026 to use the L4 \\
fp16 on the T4 corrupts the answers & W06 & High & Data-type check on 73 questions: fp16 broke 23; switched to bf16 on the L4 (0/73 broken) \\
Data from the wrong source & W05 & High & The local data was ViVQA-X rather than ViVQA, and ViTextVQA came from an unofficial copy. Re-downloaded ViVQA from the original repository, requested access to the gated ViTextVQA repository, and used the answers in \path{test_gt.json} \\
A full training run takes about 12 hours & W08 & Medium & Ran the main test on the 8,000-question subset (B3 against B2-sub); ran the full B3-full only in W10, once the subset result was clear \\
Library API changes & W06 & Low & \texttt{transformers} 5.17.0 dropped \path{warmup_ratio}; switched to \path{warmup_steps} and recorded library versions in the logs \\
GPU PaddleOCR conflicts with torch in the prototype & W13 & Medium & Ran OCR in a separate process (\path{src/prototype/ocr_worker.py}) \\
The shared adapter is biased towards ``yes'' & W12--W13 & Medium & Chose a threshold on the first-token logit using a calibration set separate from POPE-vi \\
Images with an EXIF rotation flag corrupt the OCR cache & W14 & Medium & Checked with a pre-registered rule (164/165 images with a rotation flag were misaligned); fixed \path{ocr.py}, re-ran OCR on 167 images, re-ran inference for B3-full (EM 53.35 $\rightarrow$ 53.63) \\
Review schedule changed & W10 & Low & Review~2 moved to 29 October 2026; the team updated one cumulative slide deck instead of starting over \\
\bottomrule
\end{tabular}
\end{table}

\subsection{Quality management}
\label{sec:pm-quality}

In this project, quality is first of all the reliability of the numbers. The team uses five
measures:

\begin{itemize}
  \item \textbf{Pre-registered rules:} the threshold and the way to conclude for each main test
  are committed before the results exist, for example the rule for the shared adapter (commit
  \texttt{324caeb}, W11) and the H4 hallucination rule (commit \texttt{9209170}, W12).
  \item \textbf{Paired tests:} every comparison of two models on the same test set uses a
  per-question McNemar test for EM and 95\% paired-bootstrap confidence intervals for ANLS; the
  main result is repeated with seeds 42, 43 and 44.
  \item \textbf{Recomputing from the original data:} the final result table
  (\path{src/eval/full_eval_table.py}) reads the \path{predictions.json} files directly, without
  copying numbers from the weekly reports.
  \item \textbf{Self-tests:} the main modules have a \texttt{-{}-selftest} mode that runs on the
  CPU, on both Windows and Linux, to catch bugs before spending GPU time.
  \item \textbf{Hashes and run signatures:} adapters, the OCR cache, data files and slides are
  recorded with SHA-256 hashes; each run carries a configuration signature, so the scripts do
  not append results of a different configuration to the same directory.
\end{itemize}

\paragraph{Integrating the components.} The components of the system are connected through a few fixed shared points. A
single prompt-building function is used for training, evaluation and the prototype,
so the prompt never differs between stages. OCR is run once and cached with a hash;
training and evaluation only read the cache. The list of official models is generated
automatically from the logs, and the prototype is checked against pre-set rules to
make sure it gives the same answers as the evaluation pipeline.

\subsection{Budget and compute resources}
\label{sec:pm-budget}

The initial estimate in W01 assumed a university GPU server. In W06 the team confirmed that
there was no GPU server, so costs were re-estimated in Google Colab compute units (CU). From
W08 on, before each GPU task, the team estimated the GPU hours and CU (L4 about 1.54 CU/hour,
T4 about 1.2 CU/hour) and compared them with the actual readings on Colab. A full QLoRA
training run on ViTextVQA takes about 12 L4 hours (B2: 11.99 hours; B3-full: 12.75 hours), so
the number of full training runs is the main constraint of the plan.

The Colab budget is tracked in CU and recorded in the weekly reports. The remaining CU read on
Colab were: 146.75 (13 Oct), 86.52 (19 Oct), 49.3 (21 Oct), 34.96 (23 Oct) and 17.38 (31 Oct).
W10 alone, with the LoRA configuration ablations and B3-full, used about 37~CU.

\subsection{Change management process}
\label{sec:pm-change}

Changes to the scope, the schedule, the method or the resources follow one procedure.
(1) A change is triggered by evidence --- a measurement, a reviewer's comment or a
supervisor's decision --- and never by a preference formed after seeing the results.
(2) The change, its reason and its expected impact are written in the week's README or Team
Report before it is applied. (3) Changes that affect the main results, such as the baseline,
the data source or the evaluation protocol, are confirmed with the supervisors. (4) Earlier
runs are never overwritten: superseded runs stay in \path{experiments/}, and the decision
rule of the new run is committed before the run. Table~\ref{tab:pm-changes} lists the main
changes made during the project.

\begin{table}[htbp]
\caption{Main changes during the project}
\label{tab:pm-changes}
\centering
\small
\begin{tabular}{p{0.08\textwidth}p{0.40\textwidth}p{0.42\textwidth}}
\toprule
\raggedright \textbf{Week} & \raggedright \textbf{Change} & \raggedright \textbf{Reason and handling} \tabularnewline
\midrule
\raggedright W04 & \raggedright Baseline candidate Qwen3-VL-4B replaced by Qwen2.5-VL-3B-Instruct & \raggedright Mature tooling, on the list of approved baselines, light enough for a T4; confirmed by the supervisor on 7 October 2026 \tabularnewline
\raggedright W05 & \raggedright Local ViVQA-X replaced by the original ViVQA; ViTextVQA mirror replaced by the official gated repository & \raggedright The local data turned out to be a different dataset and an unofficial copy; test answers taken from \path{test_gt.json} \tabularnewline
\raggedright W06 & \raggedright Compute moved to Google Colab (Colab Pro with an L4 from 7 October); configuration freeze moved to W07 & \raggedright No university GPU server; the T4 has no bf16, so the baseline is re-run only once, on the final hardware \tabularnewline
\raggedright W07--W08 & \raggedright OCR experiments first on a fixed 8,000-question subset; B3 on the full data only in W10 & \raggedright One ViTextVQA epoch costs about 12 L4 hours \tabularnewline
\raggedright W10 & \raggedright Hypotheses renumbered to the final H1--H4; LoRA ablations named A1 and A2 & \raggedright Review~1 feedback: the hypotheses must target Vietnamese-specific problems \tabularnewline
\raggedright W14 & \raggedright OCR cache rebuilt for images with an EXIF rotation flag & \raggedright Pre-registered rule L3: official numbers kept, the post-fix result reported as a sensitivity analysis (Section~\ref{sec:res-exif}) \tabularnewline
\raggedright W15 & \raggedright Defence demonstration moved from Colab to an offline laptop & \raggedright The demo must be shown live; the laptop starts faster (about 27 s against about 60 s) and needs no network (Section~\ref{sec:sd-deploy}) \tabularnewline
\bottomrule
\end{tabular}
\end{table}

\subsection{Closure and evaluation}
\label{sec:pm-closure}

The project is closed in three steps.

\paragraph{Freezing the results.} The six official models are listed in
\path{configs/official_models.json}, a file generated from the logs; the four adapters stored
on Google Drive match their recorded SHA-256 hashes. After this point no experiment changes
an official number: the post-fix EXIF result, for example, is reported only as a sensitivity
analysis.

\paragraph{Handover.} The reproducibility package (Appendix~\ref{appendix:repro}) lets someone
outside the team rebuild every official number and run the prototype. The submission bundle
is mapped to the four components required by the department: the thesis, the slides, the
source code with the prototype adapter, and the progress reports. The source-code archive is
checked so that it contains no real photographs and no access tokens.

\paragraph{Evaluation.} The outcome of the project is evaluated against the objectives and
completion criteria of Table~\ref{tab:objectives} and against the decision rules of the
hypotheses (Table~\ref{tab:res-hyp}); Chapter~\ref{chap:conclusion} summarises the outcome,
and Section~\ref{sec:con-reflect} records what worked well and what could have been done
better. The remaining reviews --- the department committee and the university committee ---
are handled with the same feedback procedure as Review~1 and Review~2.

%********************************** Section 2.3 ********************************
\section{Training plan}
\label{sec:pm-training}

The project requires knowledge in four different areas, so in W01 each member read at least
three papers in their area and answered 10 fixed questions for each paper, including the
mandatory question ``what can be used for the project'' (Table~\ref{tab:pm-papers}).

\begin{table}[htbp]
\caption{Foundational papers by area}
\label{tab:pm-papers}
\centering
\small
\begin{tabular}{p{0.32\textwidth}p{0.58\textwidth}}
\toprule
\textbf{Area} & \textbf{Papers read} \\
\midrule
SV1: VLM / Baseline & BLIP-2, Florence-2, Qwen2.5-VL \\
SV2: Dataset / OCR & ViVQA, ViTextVQA, TextVQA (LoRRA) \\
SV3: PEFT / Training & LoRA, QLoRA, VL-Adapter, LLaVA-1.5 \\
SV4: Evaluation / Deployment & POPE, VQAv2, Survey of Vision-Language Models (TPAMI 2024) \\
\bottomrule
\end{tabular}
\end{table}

The remaining skills were learned along the way: QLoRA training with PEFT and bitsandbytes,
PaddleOCR and VietOCR, the McNemar test, and working across a personal machine, GitHub, Colab
and Google Drive. A bilingual glossary \path{docs/Glossary.md} was kept up to date throughout
the project.

%********************************** Section 2.4 ********************************
\section{Deliverables}
\label{sec:pm-deliverables}

Table~\ref{tab:pm-deliverables} lists the deliverables required by the proposal, including the
reproducibility package. Status as of 27 November 2026.

\begin{table}[htbp]
\caption{Deliverables}
\label{tab:pm-deliverables}
\centering
\small
\begin{tabular}{p{0.38\textwidth}p{0.36\textwidth}p{0.16\textwidth}}
\toprule
\textbf{Deliverable} & \textbf{Location} & \textbf{Status} \\
\midrule
Source code: 49 Python modules (data, training, evaluation, prototype) & \path{src/} on GitHub (PLHGNAOH/ViVQA-VLM) & Done \\
30 Colab notebooks & \path{notebooks/} & Done \\
Logs of 142 runs & \path{experiments/} & Done \\
QLoRA adapters of the official models & Google Drive; hashes in \path{configs/official_models.json} & Done \\
Reproducibility package: instructions, library versions, inference script & \path{docs/REPRODUCE.md}, \path{requirements-colab.txt}, \path{src/prototype/infer.py} & Done \\
Web prototype and demo video (10 min 28 s) & \path{src/prototype/}, Google Drive & Done \\
Weekly reports W01--W15 & \path{docs/W01/} to \path{docs/W15/} & Done \\
Review~1, Review~2 and defence slides (English, Vietnamese) & \path{docs/W03/}, \path{docs/04_slides/} & Done \\
Thesis (English and Vietnamese editions) & \path{docs/03_report/overleaf/} (LaTeX source and PDF) & Done \\
\bottomrule
\end{tabular}
\end{table}

%********************************** Section 2.5 ********************************
\section{Configuration management}
\label{sec:pm-config}

\paragraph{Document management.} Documents fall into two kinds. Weekly documents are in
\path{docs/W01} to \path{docs/W15}; each directory has a README serving as the week's index
and a Team Report. Cross-cutting documents are in the numbered directories: \path{00_admin}
(conventions, checklists), \path{01_proposal} (the proposal), \path{02_references}
(literature matrix, reading notes), \path{03_report} (chapter drafts and the LaTeX source of the thesis), \path{04_slides} (review
slides). Each document has a single living copy; weekly directories link to it instead of
copying it. Old documents are moved to \path{_archive/} instead of being deleted.

\paragraph{Source-code management.} The source code is managed with Git in the public
repository \path{github.com/PLHGNAOH/ViVQA-VLM}, working on the \texttt{main} branch; the
\texttt{ai-context} branch is generated automatically by GitHub Actions to collect reports.
Data, images and adapter weights are not committed to Git (blocked by \path{.gitignore}) but
stored on Google Drive; their location and hashes are recorded in the logs and in
\path{configs/official_models.json}. Run names follow the convention in
\path{docs/00_admin/CONVENTIONS.md} and contain the week, method, model, dataset, data type,
GPU and seed, for example \path{W07_B2_qlora_r16_qwen25vl_vitextvqa_ep1_bf16_L4_seed42}. The
base model is pinned to a revision hash on Hugging Face.

\begin{table}[htbp]
\caption{Tools and infrastructure}
\label{tab:pm-tools}
\centering
\small
\begin{tabular}{p{0.26\textwidth}p{0.64\textwidth}}
\toprule
\textbf{Item} & \textbf{Tools} \\
\midrule
Languages and libraries & Python, PyTorch, HuggingFace Transformers, PEFT, bitsandbytes, \path{qwen_vl_utils} \\
OCR and Vietnamese processing & PaddleOCR, VietOCR; pyvi for word segmentation \\
Training and evaluation & Google Colab Pro: NVIDIA L4 for training and inference, Tesla T4 for OCR (W04--W06 used the free T4) \\
Personal machines & Windows 11, 4~GB GPU: development, self-tests, preprocessing, error analysis \\
Storage & Google Drive (data, adapters, large logs); GitHub (code, small logs, documents) \\
Prototype & Gradio \\
Writing code and documents & Cursor IDE \\
Slides and thesis & Canva, PowerPoint, Google Docs; LaTeX (Overleaf, then MiKTeX) \\
\bottomrule
\end{tabular}
\end{table}
